\section{Problem Formulation}
\label{sec:formulation}

\providecommand{\structurechange}[1]{#1}
\newtheorem{problem}{Problem}

\subsection{Heterogeneous Manipulator Team}
\label{sec:system_description}

Consider a team of heterogeneous robotic manipulators
\begin{equation}
 \mathcal{R}\triangleq\{r_i\}_{i=1}^{N},
 \label{eq:manipulator_team}
\end{equation}
where the team in~\eqref{eq:manipulator_team} contains $N$ manipulators and
$r_i$ has capability set
$\mathcal{C}_i$. A capability set describes the primitive actions, targets,
and embodiment conditions available to the corresponding manipulator. The
system receives a natural-language task $\gamma$ and continuous visual
observations $o(t)$ while the physical workspace evolves.

The execution process can also produce discrete semantic events
\begin{equation}
 e_k\in\mathcal{E}\triangleq
 \{e^{\texttt{t}},e^{\texttt{x}},e^{\texttt{f}}\},
 \label{eq:semantic_events}
\end{equation}
where $e^{\texttt{t}}$ denotes an online task update, $e^{\texttt{x}}$ denotes
an external disturbance, and $e^{\texttt{f}}$ denotes an execution failure.
Continuous physical motion, including motion of a target on a conveyor, remains
part of the evolving workspace and is not itself a semantic event.

\subsection{Online Primitive Command Sets}
\label{sec:online_commands}

An executable command $c$ is an instruction that can be sent directly to a
manipulator execution interface. It specifies, through the interface
representation rather than a semantic graph node, the execution skill, semantic
target or object, and manipulator that executes the skill. Thus, a semantic
subtask or graph node is not itself a command; the required output is an online
sequence of command sets
\begin{equation}
 \mathbf{C}\triangleq
 (\mathbf{c}_1,\mathbf{c}_2,\cdots,\mathbf{c}_M),
 \label{eq:command_sequence}
\end{equation}
where $M$ is determined online, each $\mathbf{c}_m$ contains capability-valid
commands that can be issued concurrently to the team, and the internal command
fields are defined in Section~\ref{sec:construct_commands}. The next command
set is generated online by the role-shared coordination
pipeline from the task, heterogeneous capabilities, visual history, semantic
event history, and previously issued command sets:
\begin{equation}
 \mathbf{c}_m=
 \Pi_{\mathrm{coord}}\!\left(
 \gamma,\{\mathcal{C}_i\}_{i=1}^{N},o(0\!:\!t_m),
 e_{1:k_m},\mathbf{c}_{1:m-1}\right),
 \label{eq:online_command_generation}
\end{equation}
where $t_m$ is the decision time and $k_m$ is the number of semantic events
observed by that time. The operator $\Pi_{\mathrm{coord}}$ in
\eqref{eq:online_command_generation} denotes the
complete online pipeline whose internal stages construct the dependency graph
$G$, executable commands, and verification status as defined
in the Method. The ordering of the arguments preserves the information flow
from task specification and embodiment capabilities to observations, events,
and prior command sets.

\subsection{Online Coordination Problem}
\label{sec:coordination_problem}

The physical workspace may change continuously between decision times, and the
event stream may contain task updates, external disturbances, or execution
failures. The command sequence must support concurrent progress while
respecting each manipulator's capability set and the dependencies induced by the
task. No scalar optimization objective is imposed in this formulation; the
primary requirement is completion of the current task under the stated online
information and execution conditions.

\begin{problem}
Given a natural-language task $\gamma$, a heterogeneous team $\mathcal{R}$ with
capability sets $\{\mathcal{C}_i\}_{i=1}^{N}$, continuous visual observations
$o(t)$, and the event history in~\eqref{eq:semantic_events}, compute the online
command-set sequence $\mathbf{C}$ in~\eqref{eq:command_sequence} such that
the primitive commands are capability-valid, independent commands can execute
concurrently, and the current task is completed despite continuous physical
evolution and the possible occurrence of $e^{\texttt{t}}$, $e^{\texttt{x}}$,
or $e^{\texttt{f}}$.\hfill $\blacksquare$
\end{problem}
